%! AP Statistics — Unit 9, Lesson 9.3 — Student Packet Cover Sheet
\documentclass[10pt]{article}
\usepackage{apstats-article}
\usepackage{apstats-boxes}
\usepackage{tikz}

\begin{document}

\begin{tikzpicture}[remember picture, overlay]
  \fill[navy] (current page.north west) rectangle
    ([yshift=-2.2cm]current page.north east);
  \node[anchor=west, white, font=\large\bfseries] at
    ([xshift=1cm, yshift=-0.7cm]current page.north west)
    {\CourseName: \SchoolYear};
  \node[anchor=west, white, font=\normalsize] at
    ([xshift=1cm, yshift=-1.3cm]current page.north west)
    {Unit 9: Inference for Quantitative Data: Slopes};
  \node[anchor=west, white, font=\normalsize\bfseries] at
    ([xshift=1cm, yshift=-1.9cm]current page.north west)
    {Lesson 9.3 \quad Justifying a Claim About the Slope of a Regression Model Based on a Confidence Interval};
\end{tikzpicture}

\vspace{2.4cm}
\namedateperiod

\begin{learningtargetbox}
\begin{itemize}
    \item I can interpret a confidence interval for the slope of a regression model in context.
    \item I can use a CI for slope to justify or refute a claim about the true slope $\beta$.
    \item I can explain how increasing sample size affects the width of a CI for slope.
\end{itemize}
\end{learningtargetbox}

\begin{tocbox}
% A CONTENTS LIST ONLY: no Score column and no Total row. Whether a component counts is said in
% its description.
\renewcommand{\arraystretch}{1.5}
\begin{tabularx}{\linewidth}{clX}
    \textbf{\#} & \textbf{Component} & \textbf{Description} \\
    \hline
    1 & Warm-Up      & Spiral review: CI for slope, claim evaluation, sample size \\
    2 & Guided Notes & CI interpretation \& claim-justification logic \\
    3 & Activity     & Three regression contexts: interpret, justify, compare \\
    4 & Exit Ticket  & Independent: interpret CI, justify claim about $\beta$ \\
    5 & Homework     & Practice interpreting CIs and evaluating researcher claims \\
\end{tabularx}
\end{tocbox}

\end{document}
